The data come from the German General Social Survey (ALLBUS) 2016, conducted by GESIS \citep{allbus2016}. To the authors' knowledge, it is the only official survey that contains both standard measures of antisemitic views and questions on Israel; the questions on Israel were included only in the 2016 wave \citep[pp.~2333, 2345]{binstok2024}.

The sample is restricted to respondents born in Germany between 1930 and 1990 who are Christian or have no religious affiliation. Jewish and Muslim respondents are excluded because they may have a personal connection to the Israeli--Palestinian conflict.
The results therefore apply to people without such a direct connection. 
After also excluding respondents with missing answers to the main variables, the sample consists of 2,336 persons \citep[pp.~2345, 2347]{binstok2024}.

The main outcome is agreement with the statement ``It is unjust that Israel is taking away land from the Palestinians'', on a scale from 1 (completely disagree) to 7 (completely agree), with a mean of 5.30. As a second measure, the authors use ``Given Israel's policies, I can well understand that people are against the Jews'', with a mean of 3.55. A weakness of the main outcome is that disagreement can mean two different things: a respondent may find the taking of land just, or may believe that no land is being taken at all \citep[pp.~2345--2346]{binstok2024}.

The main measure of antisemitism is agreement with the statement ``Jews have too much influence in the world'', on the same 1--7 scale, with a mean of 3.09. It reflects the classic conspiracy myth of a secret Jewish world power.
All items are coded so that higher values indicate more negative views of Jews.
For robustness, the authors use four further antisemitism items as well as the mean and the first factor of all five items \citep[pp.~2348--2349]{binstok2024}.

As controls, the authors use age, gender, education, regional variables (share of foreigners and unemployment rate in the county, federal state), political orientation, religion, national pride, and views on foreigners and on Islam. The authors include these controls because earlier studies show that both antisemitism and criticism of Israel are related to political views and to prejudice against other groups, such as foreigners or Muslims \citep[p.~2352]{binstok2024}. The instrument, being born between 1930 and 1945, applies to 18\% of the sample \citep[p.~2347]{binstok2024}.